\section{Related Work}
\label{sec:related}
\vspace{-1mm}

\paragraph{Multi-Armed Bandits in Education.}
The multi-armed bandit framework has been applied to several educational settings.
\citet{clement2015multi} proposed using bandits for pedagogical sequence optimization in intelligent tutoring systems, treating each learning activity as an arm.
\citet{liu2014trading} explored the exploration-exploitation trade-off in practice scheduling, showing that bandits can balance scientific knowledge discovery with student learning.
\citet{rafferty2019statistical} studied the statistical consequences of using Thompson Sampling for adaptive educational experiments and showed that bandit-based A/B tests can achieve better student outcomes than fixed designs.
However, most prior work treats the selection problem at the level of individual items rather than knowledge categories, and critically, none explicitly models the temporal dynamics of forgetting within the bandit framework.

\paragraph{Non-Stationary and Restless Bandits.}
Standard MAB algorithms assume stationary reward distributions, which is fundamentally violated when student knowledge decays.
\citet{garivier2011upper} introduced sliding-window UCB (SW-UCB) and discounted UCB for piecewise-stationary environments, providing regret bounds that depend on the number of changepoints.
The restless multi-armed bandit (RMAB) formulation~\citep{whittle1988restless} generalizes the classical bandit to settings where arm states evolve independently of the player's actions.
Whittle indices provide a tractable heuristic for the RMAB, which is PSPACE-hard in general~\citep{papadimitriou1999complexity}.
Recent work on RMAB for education includes EduQate~\citep{mate2024eduqate}, which uses restless bandits for student engagement but does not model Ebbinghaus forgetting or provide the comprehensive algorithmic comparison we present.
Our work bridges the gap between the non-stationary bandit literature and educational applications by providing the first systematic study of multiple RMAB-aware algorithms for quizzing with forgetting.

\paragraph{Knowledge Tracing and Forgetting.}
Bayesian Knowledge Tracing (BKT)~\citep{corbett1994knowledge} models student mastery as a hidden Markov model and has been widely adopted in intelligent tutoring systems.
Deep Knowledge Tracing (DKT)~\citep{piech2015deep} applies recurrent neural networks to model knowledge from interaction sequences, achieving strong predictive performance but sacrificing interpretability.
The forgetting curve, first characterized by \citet{ebbinghaus1885memory}, describes exponential decay of memory strength over time.
\citet{settles2016trainable} developed a trainable spaced repetition model for Duolingo that incorporates half-life regression to predict word forgetting.
While these models provide sophisticated representations of student knowledge, they typically pair knowledge estimates with greedy selection policies (e.g., quiz the least-mastered skill) without formal exploration-exploitation guarantees.
Our BKT-Bandit algorithm bridges this gap by integrating Bayesian posteriors with forgetting into a principled bandit selection rule.

\paragraph{Spaced Repetition Systems.}
Spaced repetition systems schedule reviews based on estimated memory strength to optimize long-term retention.
The Leitner system~\citep{leitner1972so} uses a simple box-based heuristic: correctly answered items move to higher boxes (reviewed less frequently) while incorrect items return to lower boxes.
\citet{tabibian2019enhancing} formulated spaced repetition as a stochastic optimal control problem and derived optimal review policies under a marked temporal point process model of memory.
\citet{reddy2016unbounded} proposed an RMAB-inspired approach to spaced repetition with theoretical guarantees.
These systems focus on retention of \emph{known} material rather than discovery of \emph{unknown} weaknesses, and do not incorporate the exploration component central to bandit-based approaches.
Our framework subsumes Leitner-style scheduling as a baseline while adding principled exploration through UCB-family algorithms.

\paragraph{Discounted and Non-Stationary Thompson Sampling.}
\citet{raj2017taming} proposed discounted Thompson Sampling (dTS), which geometrically discounts past observations to handle non-stationary bandits, shrinking the posterior toward the prior at each timestep.
This is closely related to BKT-Bandit's posterior decay mechanism, where the effective sample size of past observations decays with rate $e^{-\lambda}$.
We include discounted TS as a baseline and find it performs at the level of standard Thompson Sampling, suggesting that naive discounting without structural knowledge of the forgetting model provides limited benefit.
In the expert aggregation setting, EXP4~\citep{lattimore2020bandit} and Corral~\citep{agarwal2017corralling} provide formal regret guarantees for combining multiple bandit algorithms.
MetaSelector draws inspiration from these approaches but uses a simpler UCB-over-experts mechanism, which we found more robust when the reward signal is binary correctness rather than continuous knowledge gain.

Our work is distinguished from all prior approaches by simultaneously addressing three challenges: (1) principled exploration-exploitation under (2) structured non-stationarity from forgetting, with (3) a comprehensive empirical comparison of 16 algorithms across diverse regimes, including the first identification of the Greedy-Min catastrophe phenomenon.
